\documentclass[exercice,a4]{thedocument}%
\usepackage{texmaths-tikz}%
\usepackage{pstricks}%
\usepackage{pst-bezier}%
\usepackage{pstricks-add}%
\usepackage{pst-eucl}%
\usepackage{pst-func}%
\usepackage{pst-tree}%
\usepackage{pst-circ}%
\usepackage{tikz}%
\usepackage{pgfplots}%
\usepackage{tkz-tab}%
\usepackage{tikzpagenodes}%
\usepackage{tabularray}%
\usepackage{tabularx}%
\usepackage{multirow}%
\usepackage{multicol}%
\usepackage{scratch3}%
\usepackage{enumitem}%
\usepackage[np]{numprint}%
\usepackage{pas-tableur}%
\usepackage{listings}%
\usepackage{graphicx}%
\usepackage{adjustbox}%
\usepackage{needspace}%
\usepackage{xcolor}%
\usepackage{eurosym}%
\usepackage{diagbox}%
\usepackage{colortbl}%
\usepackage{multido}%
\usepackage{icomma}%
\usepackage[french]{babel}%
\startdatakeys
    \source{26GENMATAA1}
    \BO{2026}
    \cycle{4}
    \competencesDuSocle{CA,RA,CO}
    \objectifsApprentissage{C4A1aa,C4A2ad,C4B2cb,C4D2aa,C4D3ac}
    \motsCles{fonction affine,équation}
\enddatakeys
\definecolor{gris}{gray}{0.5}%
\providecommand{\wideparen}[1]{\overset{\frown}{#1}}%
\begin{document}%
\providecommand{\degres}{^\circ}%
Lola souhaite acheter un smartphone. Elle étudie deux propositions.

\begin{center}
	\begin{tblr}{width=0.9\linewidth,colspec={|X[7cm]| X |X[7cm]|},hline{1,Z}={1,3}{solid}}
		\textbf{Offre A} && \textbf{Offre B} \\
		Le client paie 175~\euro{} à l'achat, puis son abonnement est de 16~\euro{} par mois avec un engagement de 24 mois minimum. &&
		Le client ne paie rien à l'achat, puis l'abonnement est de 23~\euro{} par mois avec un engagement de 24 mois minimum.
	\end{tblr}
\end{center}

\begin{enumerate}
	\item Au bout de 24 mois, laquelle des deux offres est la plus intéressante ?
	\item $x$ est un nombre positif qui représente le nombre de mois. On exprime le prix de ces deux tarifs en fonction de $x$, avec les fonctions suivantes :

	\begin{itemize}
		\item $f(x) = 175 + 16x$
		\item $g(x) = 23x$
	\end{itemize}

	\begin{enumerate}[label={\textbf{\alph*.}}]
		\item Associer chaque fonction à l'offre correspondante (A ou B).

		\emph{Aucune justification n'est attendue.}
		\item Au bout de combien de mois paie-t-on le même prix avec ces deux offres ?
		\item Est-on encore dans la période d'engagement ?
	\end{enumerate}
\end{enumerate}

\newpage
\medskip
\showthesource
\end{document}
